\section{Experiments}\label{sec:experiments}

\paragraph{Setup.}\label{sec:setup} All models share one U-Net (\citealp{dhariwal2021adm}, as implemented in \texttt{torchcfm}; 1.08M parameters), trained with Adam \citep{kingma2015adam}, batch 128, 50 epochs and no EMA. I-CFM and OT-CFM (OT batch 128) have three training seeds, the DDPM reference (sampled with DDIM) two; each training-side variant changes one setting of the seed-0 model. Appendix~\ref{app:setup} lists all runs and their compute, Appendix~\ref{app:extra} the full analysis behind each paragraph below.

\subsection{How reliable is the evaluation? (RQ1)}\label{sec:c1}
\begin{wrapfigure}{r}{0.43\textwidth}
\vspace{-1.5em}
\centering
\includegraphics[width=0.43\textwidth]{figs/fig_samples.pdf}
\caption{MNIST samples from the same eight noise vectors after 1, 2, 5 and 100 Euler steps (DDIM uses its own grid). One I-CFM step lands on the data mean, as \eqref{eq:postmean} predicts; OT-CFM already separates digits; DDIM needs 5 steps. After reflow, one step gives the digit that 100 base steps give.}
\label{fig:samples}
\vspace{-1em}
\end{wrapfigure}
\begin{figure}[t]
\centering
\includegraphics[width=0.7\linewidth]{figs/fig_c1_reliability.pdf}
\caption{Three training seeds per model. (a) FID against Euler steps (10,000 samples). (b) Total variation between generated and real $k$-means cluster shares ($k=10$); dash-dotted: held-out real test set. Lower is better in both panels. I-CFM seed~2 (dotted) has much worse FID but the same class balance as seed~0.}
\label{fig:c1}
\end{figure}
Training seeds dominate the error (Figure~\ref{fig:c1}a; Table~\ref{tab:c1} in Appendix~\ref{app:c1}). At NFE~100 the seed sd of FID is 5.8 (I-CFM) and 1.9 (OT-CFM), while new sampling noise changes FID by an sd of only 0.13 and 0.05. Separately trained models therefore need differences above $2\hat\sigma_{\text{seed}}$ (11.6 and 3.8 at NFE~100), whereas samplers on one model can be compared to within a fraction of a point. By this rule OT-CFM beats I-CFM at 1, 2 and 5 steps on every seed (by about 110 points at one step) and cannot be distinguished from it from 10 steps on.
The course method sees what FID misses (Figure~\ref{fig:c1}b). The two-sample $\chi^2$ test does not reject the held-out real set (TV $=0.017$, $p=0.14$, near the floor 0.014), whereas a one-sample test that treats $p_k$ as known would ($p=0.037$). One-step I-CFM puts almost all samples into one cluster (TV 0.70--0.79), one-step OT-CFM does not (0.08--0.14). I-CFM seed~2 has the worst FID but the same TV as seed~0 (0.098 versus 0.092): its problem is image quality, not class balance.

\subsection{Inference side: solvers and time grids (RQ2)}\label{sec:c2}
For both models 65--69\% of $\int\kappa$ lies in $t>0.9$ (Figure~\ref{fig:c2}a in Appendix~\ref{app:c2}), and \eqref{eq:sqrtk} predicts error ratios $G=0.50$ (OT-CFM) and 0.54 (I-CFM). With Euler the $\sqrt\kappa$ grid lowers FID at 32 steps from 16.2 to 12.5 (OT-CFM) and from 18.4 to 15.2 (I-CFM), within 10\% of dopri5 at about 30 instead of 270 NFE; it helps at 16 and 32 steps on all six models (by 0.5--3.7) and is slightly worse at 2--4 steps.
\paragraph{Second-order solvers fail because of $t=1$.} On the uniform grid Heun and RK4 are far worse than Euler (OT-CFM, 32 NFE: 50.4 and 40.7 versus 16.2) because only they call the network at $t=1$. There the ideal output is $u_1(x)=x$ for independent pairs (Proposition~\ref{prop:onestep} at $t=1$), but our networks miss it by 34--55\% (RMS) on real and 83--91\% on generated images. Making only the last step avoid $t=1$ (Heun*, RK4*) takes Heun at about 31 NFE from 48.0--58.1 to 14.4--24.5 on all six models. The best samplers reduce NFE@20 from 9.2 to 6.6 (OT-CFM) and from 23.7 to 9.9 (I-CFM) at no training cost. Choosing the grid by endpoint error against dopri5 instead makes FID worse in all ten Euler cases (Appendix~\ref{app:c2}).

\subsection{Training side: pairing, time distribution and reflow (RQ3)}\label{sec:c35}
\begin{table}[t]
\centering
\caption{Training-side variants, training seed 0, Euler sampler on the uniform grid, FID on 10,000 samples (lower is better). $R$: path-length ratio (1 = straight). NFE@20: steps to reach FID 20. Time: training wall time (reflow: plus generating the pairs) relative to the base model, both throttled to the same duty cycle (OT batch 2,048 was resumed after a power loss at epoch 20; its time adds the 20 epochs before the resume). Bold: differs from the base model of the same method by more than the improvement threshold $2\,\mathrm{sd}$ over training seeds (last row).}
\label{tab:training}
\small
\setlength{\tabcolsep}{2.9pt}
\begin{tabular}{lccccccccc}
\toprule
 & \multicolumn{5}{c}{FID at NFE} & & & & \\
\cmidrule(lr){2-6}
Model & 1 & 2 & 5 & 10 & 100 & dopri5 & $R$ & NFE@20 & Time \\
\midrule
I-CFM (base) & 332 & 132 & 45.3 & 27.9 & 14.5 & 12.1 & 1.021 & 23.7 & 1.00 \\
OT-CFM (base) & 223 & 76.0 & 26.7 & 19.0 & 13.9 & 11.5 & 1.011 & 9.15 & 1.00 \\
OT-CFM, OT batch 512 & \textbf{195} & 70.4 & 30.6 & 23.5 & 17.4 & \textbf{16.7} & 1.010 & 22.1 & 1.30 \\
OT-CFM, OT batch 2,048 & 208 & 82.9 & 37.8 & 26.4 & 17.4 & \textbf{15.6} & 1.009 & 31.9 & 1.97 \\
I-CFM, logit-normal $t$ & 327 & \textbf{103} & 30.0 & 18.3 & 12.6 & 14.2 & 1.020 & 9.06 & 1.22 \\
I-CFM, $\kappa$-weighted $t$ & 335 & 122 & 37.2 & 26.0 & 18.3 & 17.5 & 1.024 & 58.1 & 1.10 \\
I-CFM + reflow & \textbf{64.6} & \textbf{30.9} & \textbf{23.9} & 22.6 & 20.5 & 19.0 & 1.001 & $>100$ & 2.67 \\
OT-CFM + reflow & \textbf{61.2} & \textbf{28.9} & 23.8 & 22.6 & \textbf{21.5} & \textbf{21.1} & 1.001 & $>100$ & 2.83 \\
\midrule
\multicolumn{10}{l}{Threshold $2\,\mathrm{sd}$ over three training seeds, I-CFM / OT-CFM:} \\
 & 18.6/19.8 & 10.2/7.81 & 15.6/11.6 & 14.1/7.58 & 11.6/3.76 & 10.1/3.41 & & & \\
\bottomrule
\end{tabular}
\end{table}
% sources: results/e4/mnist_eval.json, results/e3/fid_s0_*.json, results/e2/*_train.json, results/c5_reflow/pairs_*.json

Table~\ref{tab:training} compares each variant with its seed-0 base model; bold entries exceed the seed threshold.
\paragraph{Larger OT batches do not help.} Exact OT saves 4.5\% of the random pairing cost at batch 128 and only 6.1\% at 4,096 (Table~\ref{tab:pairs}). OT batch 512 has a lower one-step FID (196) than all three base seeds (205--223), but beyond the threshold only against seed~0, not against the seed mean; batch 2,048 has the straightest OT-CFM paths, yet is 3--11 points worse than seed~0 from two steps on and doubles the training time.
\paragraph{Where to train matters more than how to pair.} Logit-normal times, which rarely sample $t$ near 0 or 1, lower I-CFM's NFE@20 from 23.7 to 9.1, as much as the best sampler, and the model is better than all three base seeds at every NFE from 2 to 100. With a single training seed, however, only the two-step gain (29.4) exceeds the threshold. Drawing half of the times where $\kappa$ is large gives the curviest paths ($R=1.024$) and no detectable change: the sampler should spend steps near $t=1$, training should not spend samples there.
\paragraph{Reflow straightens the paths but lowers the ceiling.} Reflow pairs save 12.7--12.9\% of the pairing cost and the loss falls from 0.168 to 0.004, as Proposition~\ref{prop:cfm} predicts for a deterministic coupling ($C\approx0$). Paths become almost straight ($R\approx1.001$) and one-step FID drops from 332 to 64.6 (I-CFM) and from 223 to 61.2 (OT-CFM). With many steps both models plateau near 21 (dopri5: 19.0 and 21.1 instead of 12.1 and 11.5), beyond the threshold for OT-CFM from 20 steps on: the reflowed model can at best reproduce its base model's ODE samples, and adds the error of a second regression (Appendix~\ref{app:d7}). After reflow I-CFM and OT-CFM are nearly identical. Reflow costs 2.7--2.8 times the base training time; class balance is unharmed (TV 0.066--0.069).

\subsection{Synthesis (RQ4)}\label{sec:c6}
\begin{figure}[t]
\centering
\includegraphics[width=0.8\linewidth]{figs/fig_frontier.pdf}
\caption{Quality--efficiency frontier, training seed~0 (FID, 10,000 samples, lower is better). Solid/dash-dotted: Euler on the uniform grid; dotted: improved samplers on the $\sqrt\kappa$ grid (Heun* for base models, Euler for reflow). Grey: lower envelope. Variants without a detectable gain are omitted.}
\label{fig:frontier}
\end{figure}
The best change depends on the budget (Figure~\ref{fig:frontier}). Reflow forms the lower envelope at 1--4 NFE (61.2 at one step against 223 for the best base model); from 7 NFE on the improved samplers on OT-CFM and logit-normal I-CFM with plain Euler alternate within about one FID point (12.4 at 31 NFE, 12.0 at 50). Reflow buys a 2.6--3.6 times lower one- and two-step FID for 2.7--2.8 times the training time but never reaches FID 20; the samplers cut NFE@20 by 28--58\% for free. On the reflow models Heun and RK4 still fail on the uniform grid although the paths are nearly straight, so that failure is not caused by curvature. Over twelve models the Spearman correlation between $R$ and one-step FID is 0.92, but only 0.3 and 0.5 within OT-CFM and I-CFM: straightness separates the groups, not the models within them. Published CIFAR-10 results show the same pattern at larger scale (Table~\ref{tab:literature}): a tuned Heun sampler reaches the best FID with 35 NFE \citep{karras2022edm}, and reflow makes one step usable \citep{liu2023rf}.
